% ===== roadmap of this week =====
\begin{frame}{이번 주차 개요}
  \begin{block}{이 주차를 마치면}
    \small
    \begin{itemize}
      \item XOR이 직선 하나로 안 풀리는 이유와 은닉층의 역할
      \item 역전파를 연쇄법칙으로 유도하고 수치 미분으로 검증
      \item 그래디언트 소실·폭발 진단, 활성함수·초기화 선택
      \item 정규화와 학습률 스케줄링으로 학습 안정화
    \end{itemize}
  \end{block}
  \begin{itemize}
    \item \kb{9.1 퍼셉트론과 역전파} — XOR의 벽, 2-2-1 손계산, numpy 구현
    \item \kb{9.2 소실·폭발과 초기화} — ReLU 계열, Xavier/He 초기화
    \item \kb{9.3 정규화와 학습률} — dropout, BatchNorm, warmup·코사인
  \end{itemize}
  \vskip0.4em
  {\footnotesize 세 절은 한 사슬: 그래디언트는 ``연쇄법칙의 곱''(9.1) $\to$
  깊어지면 0으로 사라지거나 불어나는(9.2) $\to$ 학습을 안정적으로 끝내는 도구(9.3).
  Block B(CNN $\to$ RNN $\to$ 트랜스포머 $\to$ LLM)의 서막이다.}
\end{frame}
